\chapter{Protocols I: FIX}\label{ch:ll:protocols-i-fix}

An order sent this morning from a fund to its broker most likely travelled as a line of text: \texttt{35=D} for a new
order, \texttt{55=ESZ6} for the instrument, \texttt{54=1} for a buy, \texttt{38=5} for the quantity, each pair ended by a
control character. The format was first written in 1992 so that Fidelity Investments and Salomon Brothers could exchange
equity trading information electronically, and it became the common language of orders between the buy side, brokers and
many venues: the Financial Information eXchange protocol. It is old, verbose, slow to parse naively, and everywhere. This
chapter builds a \omterm{def:ll:protocols-i-fix:engine}{FIX engine}: the codec that reads and writes the text without copying it, the session layer that keeps two
firms' views of a conversation identical through disconnections, and the measurements that say what the protocol costs.

\section{Tag-value encoding and the message}

\begin{definition}[FIX protocol, tag-value encoding]\label{def:ll:protocols-i-fix:fix}
The \emph{FIX protocol}\index{FIX protocol} is a standard for messages about trading (orders, executions, quotes, market
data, allocations) between firms, maintained by the FIX Trading Community, with a session layer that delivers them in order
and without loss. In its \emph{tag-value encoding}\index{tag-value encoding} a message is a sequence of fields
\texttt{tag=value}, each ended by the byte SOH (value 1), where the tag is a number from the protocol's dictionary and the
value is text.
\end{definition}

Three fields open every message, in a fixed order: BeginString (tag 8, the protocol version), BodyLength (9) and MsgType (35).
One closes it: CheckSum (10), whose field and trailing SOH mark the end of the message. Between them come the rest of the
standard header (the sender's and target's identifiers 49 and 56, the sequence number 34, the sending time 52 in UTC) and the
body. BodyLength counts the bytes from the start of the MsgType field up to, not including, the CheckSum field, delimiters
included; CheckSum is the sum of every byte before it modulo 256, written as three digits.
\Cref{fig:ll:protocols-i-fix:anatomy} takes apart the example that the protocol's usual description works through, a logon
whose body is 65 bytes long and whose checksum is 062; the build's tests reproduce both numbers.

\begin{omfigure}
\resizebox{\linewidth}{!}{%
\begin{tikzpicture}[font=\scriptsize,f/.style={draw=gray!60,inner sep=2pt,minimum height=0.5cm,font=\ttfamily\scriptsize}]
\node[f,fill=gray!10,anchor=west] (a) at (0,0) {8=FIX.4.2|};
\node[f,fill=gray!10,anchor=west] (b) at (a.east) {9=65|};
\node[f,fill=omDef!20,anchor=west] (c) at (b.east) {35=A|};
\node[f,fill=omDef!12,anchor=west] (d) at (c.east) {49=SERVER|};
\node[f,fill=omDef!12,anchor=west] (e) at (d.east) {56=CLIENT|};
\node[f,fill=omDef!12,anchor=west] (g) at (e.east) {34=177|};
\node[f,fill=omDef!12,anchor=west] (h) at (g.east) {52=20090107-18:15:16|};
\node[f,fill=omThm!12,anchor=west] (i) at (h.east) {98=0|};
\node[f,fill=omThm!12,anchor=west] (j) at (i.east) {108=30|};
\node[f,fill=gray!10,anchor=west] (k) at (j.east) {10=062|};
\draw[decorate,decoration={brace,amplitude=4pt},thick] (c.north west) ++(0,0.08) -- ([yshift=0.08cm]j.north east)
  node[midway,above=5pt]{BodyLength: 5 + 10 + 10 + 7 + 21 + 5 + 7 = 65 bytes, from \texttt{35=} up to \texttt{10=}};
\draw[decorate,decoration={brace,amplitude=4pt,mirror},thick] ([yshift=-0.08cm]a.south west) -- ([yshift=-0.08cm]j.south east)
  node[midway,below=5pt]{CheckSum: every byte before \texttt{10=}, delimiters included, sums to 4\,158; $4\,158 \bmod 256 = 62$};
\node[anchor=north west,align=left,font=\footnotesize] at (0,-1.35) {grey: framing (first two fields, last field); blue:
  MsgType and the rest of the standard header; red: the body of a logon (encryption method 98, heartbeat interval 108)};
\end{tikzpicture}}

\omcaption{A FIX~4.2 logon as a byte string, with \texttt{|} standing for the SOH delimiter. BodyLength and CheckSum are
computed over the spans shown; a receiver recomputes both and rejects the message if either differs.}\label{fig:ll:protocols-i-fix:anatomy}
\end{omfigure}

The checksum is a weak check, and knowing how weak tells a reader what it is for.

\begin{proposition}[What the FIX checksum detects]\label{prop:ll:protocols-i-fix:checksum}
Changing a single byte of a message always changes its checksum. Swapping two bytes, or two fields, never does; nor does any
set of changes whose values sum to a multiple of 256.
\end{proposition}

\begin{proof}
Changing byte $b$ into $b' \ne b$ changes the sum by $b' - b$, with $0 < |b' - b| < 256$, which is not a multiple of 256.
A sum does not depend on the order of its terms, so a permutation leaves it unchanged; and changes $\delta_1, \dots, \delta_k$
leave the sum modulo 256 unchanged exactly when $\sum_j \delta_j \equiv 0 \pmod{256}$.
\end{proof}

The transport underneath, TCP, already carries its own checksums on every segment. The FIX checksum and body length are
therefore less a defence against corruption on the wire than a check on the framing: a program that has lost its place in a
byte stream, or a message cut short or stitched together by a faulty gateway, is caught at the next message. Framing itself
uses the first two fields: once \texttt{8=} and \texttt{9=} are read, the message's full length is known, the body plus the
seven bytes of \texttt{10=nnn} and its delimiter.

\section{The session layer: sequence numbers, heartbeats, resend and gap fill}

\begin{definition}[FIX session, heartbeat, resend request, gap fill]\label{def:ll:protocols-i-fix:session}
A \emph{FIX session}\index{FIX session} is the bidirectional conversation between two parties identified by their sender and
target identifiers, from a logon to a logout, in which each side numbers the messages it sends 1, 2, 3, \dots\ and keeps
them so that it can send them again. A \emph{heartbeat}\index{heartbeat} is a message sent when a side has sent nothing for
an agreed interval, so that silence can be told from a dead connection. A \emph{resend request}\index{resend request} asks the
other side to send again a range of its messages, after a gap in their sequence numbers has been detected. A \emph{gap
fill}\index{gap fill} is a sequence-reset message, sent in answer to a resend request, that stands for a run of
administrative messages that will not be sent again, moving the receiver's expected sequence number past them.
\end{definition}

Sequence numbers (One Quant Book~1, chapter~28, for market data) are the session layer's whole mechanism. Each side keeps two
counters: the number of its next outgoing message and the number it expects next from the other side. A message with the
expected number is processed and the counter moves on. A higher number means messages were lost: the receiver sends a
ResendRequest (MsgType 2) from the expected number to ``infinity'' (EndSeqNo 0), and holds the later messages until the gap is
filled. A lower number is either a duplicate, marked with PossDupFlag (43) and ignored, or a fatal error: the two sides no
longer agree, and the receiver logs out. The answer to a \omterm{def:ll:protocols-i-fix:session}{resend request} contains the application messages again, marked
PossDupFlag with their original sending time (122), and replaces administrative messages (\omterm{def:ll:protocols-i-fix:session}{heartbeats}, test requests, \omterm{def:ll:protocols-i-fix:session}{resend
requests} themselves) by SequenceReset messages with GapFillFlag (123) set: a \omterm{def:ll:protocols-i-fix:session}{heartbeat} from half an hour ago means nothing
now. \Cref{fig:ll:protocols-i-fix:sequence} draws the golden conversation of the build around its gap, and
\cref{lst:ll:protocols-i-fix:trace} is the same conversation as the engine traces it.

\begin{omfigure}
\begin{tikzpicture}[font=\scriptsize,>=Stealth]
\def\L{0}\def\R{7.5}
\node[font=\footnotesize\bfseries] at (\L,0.45) {FIRM (us)};
\node[font=\footnotesize\bfseries] at (\R,0.45) {BROKER};
\draw[gray,thick] (\L,0.15) -- (\L,-9.6);
\draw[gray,thick] (\R,0.15) -- (\R,-9.6);
\newcommand{\msgR}[3]{\draw[->,#3] (\L,#1) -- node[above,sloped]{#2} (\R,#1-0.25);}
\newcommand{\msgL}[3]{\draw[<-,#3] (\L,#1-0.25) -- node[above,sloped]{#2} (\R,#1);}
\msgR{-0.1}{Logon, seq 1}{omDef}
\msgL{-0.75}{Logon, seq 1}{omDef}
\msgR{-1.4}{NewOrderSingle, seq 2}{omDef}
\msgL{-2.05}{ExecutionReport (new), seq 2}{omDef}
\msgL{-2.7}{Heartbeat, seq 3}{gray}
\draw[omThm,dashed,->] (\R,-3.3) -- (5.6,-3.45) node[left,font=\scriptsize]{lost: seq 4 (partial fill), seq 5 (heartbeat)};
\msgL{-4.0}{ExecutionReport (filled), seq 6}{omDef}
\node[anchor=east,omThm,align=right] at (-0.15,-4.25) {gap: expected 4,\\ hold seq 6};
\msgR{-4.65}{ResendRequest 7=4 16=0, seq 3}{omThm,thick}
\msgL{-5.3}{ExecutionReport (partial), seq 4, PossDup}{omDef}
\msgL{-5.95}{SequenceReset-GapFill, seq 5, NewSeqNo 6}{omMeth}
\node[anchor=east,omMeth,align=right] at (-0.15,-6.2) {deliver 4, then the\\ held 6; expect 7};
\msgL{-6.6}{ExecutionReport, seq 6 again, PossDup: ignored}{gray}
\msgR{-7.25}{OrderCancelRequest, seq 4}{omDef}
\msgL{-7.9}{ResendRequest 7=3 16=0 (the broker lost ours)}{omThm}
\msgR{-8.55}{GapFill seq 3 (our resend request), NewSeqNo 4}{omMeth}
\msgR{-9.2}{OrderCancelRequest again, seq 4, PossDup}{omDef}
\end{tikzpicture}

\omcaption{The build's golden conversation around its gap. Two of the broker's messages are lost; the next one reveals the
gap, which a \omterm{def:ll:protocols-i-fix:session}{resend request} closes: the application message comes again marked as a possible duplicate, the \omterm{def:ll:protocols-i-fix:session}{heartbeat} is
replaced by a \omterm{def:ll:protocols-i-fix:session}{gap fill}, and the message held back is delivered in order. Later the broker asks for ours, and our
administrative message is gap-filled in turn.}\label{fig:ll:protocols-i-fix:sequence}
\end{omfigure}

\omcode{code/low-latency/15-protocols-i-fix/python/trace_short.txt}{7}{22}{The golden conversation around the gap, as
\texttt{ll\_fix.short} prints the engine's trace (the full trace is \texttt{expected\_trace.txt} in the build's data).}\label{lst:ll:protocols-i-fix:trace}

\omterm{def:ll:protocols-i-fix:session}{Heartbeats} make silence meaningful. At logon both sides agree an interval, HeartBtInt (108), here 30 seconds. A side that has
sent nothing for that long sends a \omterm{def:ll:protocols-i-fix:session}{Heartbeat} (0). A side that has heard nothing for somewhat longer (the build waits 1.2
intervals) sends a TestRequest (1) with an identifier that the answering \omterm{def:ll:protocols-i-fix:session}{heartbeat} must carry, and if that answer does not
come within another interval it declares the connection dead. The detection time is the price of a long interval: with 30
seconds, a silently dead connection is noticed after about 66, during which the counterparty's messages go nowhere and must be
resent.

\omcode{code/firm/fixengine/cpp/firm_fixengine.hpp}{267}{289}{The C++ engine's sequence check: a gap asks for a resend and
holds the message; a duplicate is ignored; a number too low ends the session.}\label{lst:ll:protocols-i-fix:seq}

\section{The application layer: orders and execution reports}

On top of the session, each application message carries one business event. A NewOrderSingle (35=D) carries the firm's order
identifier ClOrdID (11), the instrument (55), the side (54), the quantity (38), the order type (40) and, for a limit order,
the price (44); an OrderCancelRequest (35=F) names the order to cancel by its original identifier (41) and gives the request
its own new one. The answer to each is an execution report (35=8), which the specification uses ``to confirm the receipt of
an order'', to confirm changes, to relay status and fills and to reject orders; One Quant Book~10, chapter~26, describes it
from the venue's side. Its OrdStatus (39) walks the order lifecycle of One Quant Book~7, chapter~17: \texttt{0} new,
\texttt{1} partially filled, \texttt{2} filled, \texttt{4} cancelled, \texttt{8} rejected; LastQty (32) and LastPx (31)
describe the latest fill, CumQty (14) and LeavesQty (151) the totals. A partial fill of 2 of 5 contracts at 5723.25 therefore
reads \texttt{39=1}, \texttt{32=2}, \texttt{31=5723.25}, \texttt{14=2}, \texttt{151=3}.

The session and application layers meet in one place: resent messages. An execution report that arrives with PossDupFlag
may or may not have been processed before (the original might have arrived and only its successor been lost). The application
must therefore make processing idempotent, recognising a fill it has already booked by its execution identifier (17), or it
books a fill twice. The session layer delivers each sequence number once; it cannot know what the application did with
the original.

\section{Parsing fast}

\begin{definition}[FIX engine]\label{def:ll:protocols-i-fix:engine}
A \emph{FIX engine}\index{FIX engine} is the software that implements the \omterm{def:ll:protocols-i-fix:session}{FIX session} layer for an application: it frames
and validates incoming messages, maintains sequence numbers, \omterm{def:ll:protocols-i-fix:session}{heartbeats} and the store of sent messages, answers \omterm{def:ll:protocols-i-fix:session}{resend
requests}, and hands application messages to the program in order.
\end{definition}

A naive parser splits the message on SOH, converts each tag to an integer and copies each value into a map from tag to
string: several allocations per field (chapter~6), a tree insertion per field, and one pass per byte. The build's parser does
none of that. It finds every delimiter of the message with the vector scan of chapter~14, parses each tag in place, and
records for each field only its tag, offset and length in a fixed array: the values stay where they are, in the receive
buffer, and are read as string views. Body length and checksum are verified in the same pass.

\omcode{code/firm/fixengine/cpp/firm_fixengine.hpp}{40}{73}{The zero-copy parse: delimiters from the vector scan, tags parsed
in place, then the order of the framing fields, the body length and the checksum.}\label{lst:ll:protocols-i-fix:parse}

\Cref{fig:ll:protocols-i-fix:parse} measures both on execution reports of about 205 bytes. The zero-copy view parses one in
about \qty{80}{\nano\second}, some twelve million a second on one core; the map-based parser takes about
\qty{900}{\nano\second}, eleven times longer; the Python reference about \qty{5}{\micro\second}. Building is symmetric: the
fixed-buffer builder, which writes digits itself and computes the length and checksum in place, produces a new order in about
\qty{60}{\nano\second}, against about \qty{330}{\nano\second} for the same message assembled by string concatenation.

\begin{omfigure}
\begin{tikzpicture}
\begin{axis}[width=11.5cm,height=5.6cm,ybar,bar width=14pt,ymode=log,log origin y=infty,ymin=20,ymax=20000,
  symbolic x coords={view,map,python,builder,concat},xtick=data,
  xticklabels={{parse:\\ zero-copy view},{parse:\\ std::map},{parse:\\ Python reference},{build:\\ fixed buffer},{build:\\ concatenation}},
  xticklabel style={align=center,font=\footnotesize},enlarge x limits=0.12,
  ylabel={ns per message (log scale)},tick label style={font=\small},label style={font=\small},grid=major,
  grid style={gray!25},nodes near coords,every node near coord/.append style={font=\scriptsize},
  point meta=explicit symbolic,table/col sep=comma]
\addplot[fill=omDef!60,draw=omDef] table[x=key,y=ns_per_msg,meta=ns_per_msg]{figdata/low-latency/15-protocols-i-fix/measured_fix.csv};
\end{axis}
\end{tikzpicture}

\omcaption{Parsing and building FIX execution reports of about 205 bytes (1\,000 distinct messages, median of 40 passes).
Measured on a laptop (Intel Core Ultra 7 155H) under WSL2, no isolated cores; C++ with GCC~11 at \texttt{-O2}, Python
3.10. Data: \texttt{bench\_fix.py}.}\label{fig:ll:protocols-i-fix:parse}
\end{omfigure}

Parsing is rarely the whole cost of a FIX connection: the session layer writes each sent message to its store, often on disk,
before or after sending it, and the counterparty's engine, network and matching take far longer than
\qty{80}{\nano\second}. But on a gateway that carries a broker's whole flow, or on a market maker's order path through a venue
that speaks only FIX, the difference between \qty{80}{\nano\second} and a microsecond per message is the difference between a
component that disappears from the \omterm{def:ll:where-latency-comes-from:budget}{latency budget} of chapter~1 and one that owns part of it.

\section{Where FIX is still used}

FIX is the language of order routing between firms: from asset managers' order management systems to brokers, from brokers'
algorithms to venues, and for direct market access (One Quant Book~1, chapter~4), where a client's orders reach the venue
through the broker's infrastructure. Some venues accept it for order entry and drop copies, as the dated box below shows for one of them. Where speed is the
product, venues
have moved to binary protocols with fixed-offset fields (chapter~16), and the FIX Trading Community itself has standardised
binary alternatives: Simple Binary Encoding for the messages, and FIXP, a ``lightweight point-to-point protocol'' for the
session, whose own introduction names as its basis the \omterm{def:ll:protocols-i-fix:session}{FIX session} layer and exchange protocols such as SoupBinTCP. The
tag-value engine remains the one every trading firm needs somewhere.

\begin{dated}{2026-09}{FIX order entry today}\label{dat:ll:protocols-i-fix:today}
Coinbase Exchange documents three ways to trade (consulted September 2026): a REST interface ``for lower-frequency
trading'', a FIX order-entry interface ``for higher-frequency trading'', and FIX and WebSocket market data. Its FIX baseline
is FIX 5.0 SP2; sessions are logged out every Saturday; and ``\omterm{def:ll:protocols-i-fix:session}{resend requests} are not supported'': each connection starts a
new session with new sequence numbers, so a client that loses its connection recovers the state of its orders by other
means, such as the venue's drop-copy sessions. The FIXP session standard reached version 1.0 in April 2021.
\end{dated}

\section{Tutorial: a FIX engine in three languages}\label{tut:ll:protocols-i-fix:tut}

\begin{tutorial}
\textbf{Goal.} Parse and build FIX messages without copying them, replay a conversation with a gap through the session layer
in Python and C++, and measure the parsers. \textbf{End state:} \cref{fig:ll:protocols-i-fix:parse}, the trace of
\cref{lst:ll:protocols-i-fix:trace}, and green tests in Python, C++ and Rust.
\begin{tutsteps}
\item \textbf{Check the published example.} The tests decode the logon of \cref{fig:ll:protocols-i-fix:anatomy} and find a
body of 65 bytes and a checksum of 062, then change one byte (checksum error) and the declared length (body-length error).
\item \textbf{Replay the golden conversation.} \texttt{make\_fix\_fixtures.py} scripts the broker's side (a lost fill and a
lost \omterm{def:ll:protocols-i-fix:session}{heartbeat}, a \omterm{def:ll:protocols-i-fix:session}{resend request} from the broker, a silent broker); \texttt{firm\_fixengine.replay} runs our session through
it and writes the trace; the C++ test runs the C++ session through the same conversation and compares its trace with the
Python one, line for line; the Rust tests parse and rebuild every message of it byte for byte.
\item \textbf{Shorten the trace} with \texttt{python ll\_fix.py}, which prints \cref{lst:ll:protocols-i-fix:trace}.
\item \textbf{Measure} with \texttt{python bench\_fix.py}: the C++ parsers and builders on 1\,000 execution reports, and the
Python reference on the same messages.
\end{tutsteps}
\textbf{What to change next.} Delete the \omterm{def:ll:protocols-i-fix:session}{gap fill} from the broker's answer in the conversation and replay it: the session
waits for message 5 forever. Then make the session answer a message whose tag is not a number with a Reject (35=3).
\end{tutorial}

\section{Build: the FIX engine}\label{bld:ll:protocols-i-fix:fixengine}

\begin{build}
\bfield{Purpose} The firm's FIX connections: broker and venue order entry where FIX is the protocol, drop copies, and the
test counterparties of chapter~25. The order gateway of chapter~21 uses its codec for FIX venues.
\bfield{Interface} Python reference \texttt{firm\_fixengine}: \texttt{checksum}, \texttt{encode}, \texttt{decode},
\texttt{frame}, \texttt{Session} (\texttt{logon}, \texttt{send}, \texttt{on\_bytes}, \texttt{on\_timer}, \texttt{logout},
\texttt{trace}), \texttt{replay}. C++20 \texttt{firm::fix}: \texttt{View::parse}, \texttt{get}, \texttt{frame},
\texttt{Builder} (\texttt{start}, \texttt{add}, \texttt{finish}), \texttt{Session} with the same methods and trace. Rust
\texttt{firm\_fixengine}: \texttt{View::parse}, \texttt{get}, \texttt{frame}, \texttt{encode}, \texttt{checksum}.
\bfield{Rules} The parser copies nothing and rejects a message whose framing fields, body length or checksum are wrong; the
builder allocates nothing; the session numbers every message, stores every one it sends, answers \omterm{def:ll:protocols-i-fix:session}{resend requests} with
application messages marked PossDupFlag and \omterm{def:ll:protocols-i-fix:session}{gap fills} for administrative ones, holds messages received ahead of a gap, and
ends the session on a sequence number too low without PossDupFlag.
\bfield{Acceptance tests} \texttt{code/firm/fixengine/}: the published example's body length and checksum; checksum,
body-length and field-order corruptions caught; framing of concatenated and partial streams; the golden conversation's trace
reproduced by the Python reference and, line for line, by the C++ session; every message of that trace parsed and rebuilt byte
for byte in Rust; \omterm{def:ll:protocols-i-fix:session}{heartbeat}, test request and timeout; logout on a sequence number too low; reset mode.
\bfield{Stretch} A persistent message store (chapter~23's binary log), a Reject for malformed messages, FIXT~1.1 sessions
carrying FIX~5.0 application messages, and the session layer in Rust.
\end{build}

\begin{omsources}
\item FIX Trading Community, FIX Latest (Extension Pack 312) message and field dictionary, as published by the Orchimate
browser (fields 7, 8, 9, 10, 16, 34, 35, 36, 43, 49, 52, 56, 108, 112, 123; messages 0, 1, 2, 4, 5, 8, A, D, 9).
\item FIX Trading Community, \emph{FIX Performance Session Layer (FIXP)} specification (GitHub, versions 1.0 and 1.1).
\item Wikipedia, ``Financial Information eXchange'' (history; the worked body-length and checksum example).
\item Coinbase Developer Documentation, Exchange FIX API: connectivity.
\end{omsources}

\section{Exercises}

\begin{exercise}[$\star$]\label{exo:ll:protocols-i-fix:1}
Compute the body length and checksum of the \omterm{def:ll:protocols-i-fix:session}{heartbeat} \texttt{8=FIX.4.4|9=?|35=0|49=A|56=B|34=2|52=20260925-14:30:00.000|10=?|},
where \texttt{|} is SOH, and check your answer with \texttt{firm\_fixengine.encode}.
\end{exercise}

\begin{exercise}[$\star$]\label{exo:ll:protocols-i-fix:2}
With a \omterm{def:ll:protocols-i-fix:session}{heartbeat} interval of 30 seconds and the build's rules, the broker falls silent at 10:00:00. When does our session send a
test request, and when does it declare the connection dead if nothing comes back?
\end{exercise}

\begin{exercise}[$\star$]\label{exo:ll:protocols-i-fix:3}
Our session expects sequence number 52 and receives 57. What does it send, what does it do with message 57, and what must the
broker send back?
\end{exercise}

\begin{exercise}[$\star\star$]\label{exo:ll:protocols-i-fix:4}
Give two different, well-formed messages of the same length with the same checksum, and explain with
\cref{prop:ll:protocols-i-fix:checksum} why the protocol accepts the risk.
\end{exercise}

\begin{exercise}[$\star\star$]\label{exo:ll:protocols-i-fix:5}
Why are \omterm{def:ll:protocols-i-fix:session}{heartbeats} and test requests gap-filled instead of sent again? What would go wrong if they were resent?
\end{exercise}

\begin{exercise}[$\star\star$]\label{exo:ll:protocols-i-fix:6}
A gateway receives 200\,000 execution reports a second at peak. Using the measured costs, how much of a core does parsing take
with the zero-copy view, the map-based parser and the Python reference?
\end{exercise}

\begin{exercise}[$\star\star\star$]\label{exo:ll:protocols-i-fix:7}
\emph{Coding.} Make the Python and C++ sessions answer a message whose tag is not a number with a Reject (35=3) carrying
RefSeqNum (45), without advancing past it, and add the case to the golden conversation so that both traces still agree.
\end{exercise}

\begin{exercise}[$\star\star\star$]\label{exo:ll:protocols-i-fix:8}
\emph{Find the flaw.} ``To keep things simple our engine resets both sequence numbers to 1 at every reconnection, so it never
needs resend logic.''
\end{exercise}

\section{Problem: The Weekend Gap}

\begin{problem}[{Weekend problem --- resynchronising after a disconnection}]\label{pb:ll:protocols-i-fix:1}
A broker's engine keeps writing to its store the messages of a session whose connection to us is down, and resends them when
we reconnect. Assume messages of 205 bytes, a 1~Gb/s link with a round trip of \qty{1}{\milli\second}, a \omterm{def:ll:protocols-i-fix:session}{heartbeat} interval of
30 seconds, and the measured parsing costs (\texttt{measured\_fix.csv}). Consider two disconnections: a busy one, a link that
dies silently while the broker sends 200 execution reports a second, and a quiet one, 48 hours over a weekend with one
application message an hour and a \omterm{def:ll:protocols-i-fix:session}{heartbeat} whenever 30 seconds pass without one.

\medskip\noindent\textbf{Part I --- The busy outage.}
\begin{enumerate}
\item How long does the build's session take to notice the dead link, and why?
\item How many messages must the broker resend (\texttt{ll\_fix.sender\_stream})? How many are \omterm{def:ll:protocols-i-fix:session}{heartbeats}?
\item How many bytes is that, and how long do they take on the wire?
\item How long does resynchronising take with the C++ engine, and with the Python reference?
\end{enumerate}

\medskip\noindent\textbf{Part II --- The quiet weekend.}
\begin{enumerate}[resume]
\item How many messages sit in the broker's store for us after 48 hours? How many are application messages?
\item How many messages does the answer to our \omterm{def:ll:protocols-i-fix:session}{resend request} contain, with \omterm{def:ll:protocols-i-fix:session}{gap fill} and without?
\item How long does resynchronising take with the Python reference, each way?
\item Why does \omterm{def:ll:protocols-i-fix:session}{gap fill} change nothing in the busy case?
\end{enumerate}

\medskip\noindent\textbf{Part III --- The design.}
\begin{enumerate}[resume]
\item Why may administrative messages never be resent as they were?
\item What must the application do with a resent execution report marked PossDupFlag?
\item A venue that supports no \omterm{def:ll:protocols-i-fix:session}{resend requests} starts a new session at each connection. How do you learn about fills that
happened while you were disconnected?
\item What must the engine persist to survive its own restart without a gap?
\end{enumerate}

\medskip\noindent\textbf{Part IV --- The verdict.}
\begin{enumerate}[resume]
\item State the \emph{named result}: messages resent and time to resynchronise in both cases, with and without \omterm{def:ll:protocols-i-fix:session}{gap fill}.
\item Would a 1-second \omterm{def:ll:protocols-i-fix:session}{heartbeat} interval help the busy case? At what cost?
\item What if the broker's store kept only its last 10\,000 messages?
\item Should the strategy act on resent execution reports as if they had just happened?
\item How would you test the resend logic before connecting to a broker?
\item Can the Python reference keep up with the busy session in normal running?
\item Which of the two cases would you design the engine for?
\item In one sentence: what is the session layer for?
\end{enumerate}
\end{problem}

\section{Interview questions}

\begin{interviewq}[$\star$ \iqroles{developer}]\label{iq:ll:protocols-i-fix:1}
Describe the structure of a FIX message. How does a receiver know where a message ends?
\end{interviewq}

\begin{interviewq}[$\star\star$ \iqroles{developer}]\label{iq:ll:protocols-i-fix:2}
Walk through what happens when a \omterm{def:ll:protocols-i-fix:session}{FIX session} detects a gap in incoming sequence numbers.
\end{interviewq}

\begin{interviewq}[$\star\star$ \iqroles{developer}]\label{iq:ll:protocols-i-fix:3}
How would you write a FIX parser that does not allocate? What does it return?
\end{interviewq}

\begin{interviewq}[$\star\star$ \iqroles{developer}]\label{iq:ll:protocols-i-fix:4}
What are \omterm{def:ll:protocols-i-fix:session}{heartbeats} and test requests for, and how would you choose the \omterm{def:ll:protocols-i-fix:session}{heartbeat} interval?
\end{interviewq}

\begin{interviewq}[$\star\star$ \iqroles{developer}]\label{iq:ll:protocols-i-fix:5}
Why do many venues use binary protocols for order entry instead of FIX tag-value? What do they give up?
\end{interviewq}

\begin{interviewq}[$\star\star\star$ \iqroles{developer}]\label{iq:ll:protocols-i-fix:6}
Your engine restarted in the middle of the day and the broker says your sequence number is too low. What happened, and how do
you recover without losing or duplicating fills?
\end{interviewq}
